\section{Hard sync 的精确几何与稳定区间}
\label{dyn:sec:hard-sync-geometry}

本节固定一个 prefix，省略条件变量。沿用 $q_k$ 表示第 $k$ 轮冻结 anchor，$t$ 表示固定 teacher。词表 $\mathcal V$ 有限，$q_0(v)>0$、$t(v)>0$ 对所有 token 成立。假设 inner 优化确实得到 Gibbs target，随后将这个 target 完整复制为下一轮 anchor。这个假设同时排除了有限 inner 误差、模型表达能力不足和参数同步与分布同步之间的差异。

\begin{definition}[精确 Gibbs hard sync]
对任意 $\eta>0$ 和严格正分布 $q$，定义
\begin{equation*}
 T_\eta(q)(v)
 =\frac{q(v)^{1-\eta}t(v)^\eta}{Z_\eta(q)},
 \qquad
 Z_\eta(q)=\sum_{w\in\mathcal V}q(w)^{1-\eta}t(w)^\eta.
\end{equation*}
精确 hard sync 是 $q_{k+1}=T_{\eta_k}(q_k)$。$\eta_k$ 是 mirror 更新的步长；它与优化器学习率及每轮 inner steps 是不同量。
\end{definition}

严格正性使负指数也有定义。因此本节先分析所有 $\eta>0$ 的数学更新，再区分通常采用的 $0<\eta\leq1$ 区间。若初始分布有零概率，$0<\eta<1$ 的更新会一直保留这些零概率，不能恢复 teacher 的完整支持集；$\eta>1$ 的负幂则在这些 token 上失去定义。有限 logits 的 softmax 满足正性，而显式截断后的采样分布不能直接代入本节结论。

\subsection{归一化常数消去后的精确收缩}

\begin{definition}[Log-ratio 跨度与 Hilbert 距离]
对 token 函数 $f$，定义
\begin{equation*}
 \operatorname{osc}_v f_v=\max_v f_v-\min_v f_v.
\end{equation*}
对两个严格正概率分布，定义
\begin{equation*}
 d_H(q,\widetilde q)
 =\operatorname{osc}_v\log\frac{q(v)}{\widetilde q(v)},
 \qquad b_k=d_H(q_k,t).
\end{equation*}
$d_H$ 是正概率单纯形上的 Hilbert 距离。它忽略 log-density 中与 token 无关的加性常数，恰好对应概率归一化所消除的自由度。
\end{definition}

这个定义确实给出距离。非负性来自最大值减最小值；对称性来自 $\operatorname{osc}(-f)=\operatorname{osc}(f)$；三角不等式来自 $\operatorname{osc}(f+g)\leq\operatorname{osc}(f)+\operatorname{osc}(g)$。若 $d_H(q,\widetilde q)=0$，则两个分布的概率比是常数，归一化迫使这个常数等于一。

\begin{theorem}[两条 hard-sync 轨迹的精确收缩]
\label{dyn:thm:hs-pairwise-contraction}
相同 teacher 和相同步长下，任意严格正分布满足
\begin{equation}\label{dyn:eq:hs-pairwise-contraction}
 d_H(T_\eta(q),T_\eta(\widetilde q))
 =|1-\eta|d_H(q,\widetilde q).
\end{equation}
特别地，$t$ 是不动点，而且
\begin{equation*}
 b_{k+1}=|1-\eta_k|b_k.
\end{equation*}
\end{theorem}

\begin{proof}
先逐 token 相除，再取对数：
\begin{equation*}
 \log\frac{T_\eta(q)(v)}{T_\eta(\widetilde q)(v)}
 =(1-\eta)\log\frac{q(v)}{\widetilde q(v)}
 -\log Z_\eta(q)+\log Z_\eta(\widetilde q).
\end{equation*}
最后两个项与 token 无关，在跨度中抵消。对任意实数 $c$，$\operatorname{osc}(cf)=|c|\operatorname{osc}(f)$，从而得到式~\eqref{dyn:eq:hs-pairwise-contraction}。将 $\widetilde q=t$ 代入，并使用 $T_\eta(t)=t$，得到第二个结论。
\end{proof}

\begin{theorem}[常步长与变步长的精确解]
\label{dyn:thm:hard-sync}
记
\begin{equation*}
 \chi_0=1,\qquad \chi_k=\prod_{j=0}^{k-1}(1-\eta_j),
 \qquad u_v=\log\frac{q_0(v)}{t(v)}.
\end{equation*}
任意正步长序列的精确解为
\begin{equation}\label{dyn:eq:hard-closed}
 q_k(v)=\frac{t(v)e^{\chi_k u_v}}
 {\sum_wt(w)e^{\chi_k u_w}},
 \qquad b_k=|\chi_k|b_0.
\end{equation}
常步长时 $\chi_k=(1-\eta)^k$，也可以写成
\begin{equation*}
 q_k(v)=\frac{q_0(v)^{(1-\eta)^k}
 t(v)^{1-(1-\eta)^k}}
 {\sum_wq_0(w)^{(1-\eta)^k}
 t(w)^{1-(1-\eta)^k}}.
\end{equation*}
\end{theorem}

\begin{proof}
$k=0$ 时 $t(v)e^{u_v}=q_0(v)$，配分函数等于一。假设第 $k$ 轮公式成立，其归一化常数记为 $C_k$。下一轮未归一化权重为
\begin{align*}
 q_k(v)^{1-\eta_k}t(v)^{\eta_k}
 &=C_k^{-(1-\eta_k)}t(v)
       e^{(1-\eta_k)\chi_k u_v}\\
 &=C_k^{-(1-\eta_k)}t(v)e^{\chi_{k+1}u_v}.
\end{align*}
第一项不依赖 token，归一化后消失。由归纳得到分布闭式。取相对 teacher 的 log-ratio 后，归一化常数同样在跨度中抵消，得到 $b_k=|\chi_k|b_0$。
\end{proof}

\begin{corollary}[常步长的完整稳定分类]
对非 teacher 初值，精确 hard sync 的行为如下。
\begin{enumerate}
 \item $0<\eta<1$ 时，$\chi_k>0$ 且趋于零，策略从同一侧趋向 teacher。
 \item $\eta=1$ 时，第一轮更新等于 teacher，此后保持不变。
 \item $1<\eta<2$ 时，$\chi_k$ 正负交替并趋于零，策略在 teacher 两侧交替收敛。
 \item $\eta=2$ 时，$\chi_k=(-1)^k$，形成精确二周期。
 \item $\eta>2$ 时，$|\chi_k|$ 发散，log-ratio 跨度发散，策略不能趋向 teacher。
\end{enumerate}
因此完整的全局渐近稳定区间是 $0<\eta<2$。teacher KL 的每步单调下降需要另外分析。
\end{corollary}

\begin{proof}
前四项直接由式~\eqref{dyn:eq:hard-closed} 得到。第五项中，偶数轮的 $\chi_k$ 趋于正无穷，奇数轮趋于负无穷。偶数轮的概率质量集中到 $u_v$ 最大的 token 集合，奇数轮集中到最小的 token 集合；各集合内部的极限概率与 $t$ 成比例。非 teacher 初值使 $u$ 非常数，所以两个集合互不相交。

有限正 teacher 下仍有
\begin{equation*}
 \KL(q_k\Vert t)\leq\log\frac1{\min_vt(v)}.
\end{equation*}
因此这里发散的是 log-ratio 跨度和$\KL(t\Vert q_k)$，而$\KL(q_k\Vert t)$ 保持有界。$\KL(t\Vert q_k)$ 发散，是因为上述集中使某些具有正 teacher 质量的 token 概率趋于零。
\end{proof}

\begin{proposition}[变步长的充要条件]
非 teacher 初值下，$q_k\to t$ 当且仅当 $\chi_k\to0$。若每个 $\eta_k\in(0,2)$，且没有任何一步等于一，则这个条件等价于
\begin{equation*}
 \sum_{k=0}^{\infty}\eta_k(2-\eta_k)=\infty.
\end{equation*}
若某一步等于一，则此后 $\chi_k=0$，直接收敛。
\end{proposition}

\begin{proof}
$\chi_k\to0$ 时由有限词表闭式得到逐 token 收敛。反过来，概率趋向严格正 teacher 时 $b_k\to0$；因为 $b_0>0$，故 $|\chi_k|\to0$。另记单步绝对乘子 $\upsilon_k=|1-\eta_k|\in(0,1)$，则 $|\chi_K|=\prod_{k<K}\upsilon_k$。该无限乘积为零，当且仅当 $\sum_k(1-\upsilon_k)$ 发散。判据由 $-\log d\geq1-d$，以及 $d$ 接近一时 $-\log d\leq2(1-d)$ 推出。最后使用
\begin{equation*}
 1-\upsilon_k\leq1-\upsilon_k^2\leq2(1-\upsilon_k),
 \qquad 1-\upsilon_k^2=\eta_k(2-\eta_k),
\end{equation*}
得到等价条件。
\end{proof}

例如 $\eta_k=1/(k+2)$ 给出 $\chi_K=1/(K+1)\to0$。若改成 $\eta_k=1/(k+2)^2$，则
\begin{equation*}
 \chi_K=\prod_{n=2}^{K+1}\frac{(n-1)(n+1)}{n^2}
 =\frac{K+2}{2(K+1)}\longrightarrow\frac12.
\end{equation*}
第二种步长虽然每轮都执行正确方向的 hard sync，但总更新量有限，最终仍保留初始 log-ratio 的一半。

\subsection{配分函数与两个方向的 KL 路径}

沿用上节 $u_v$，定义一条由初始策略和 teacher 唯一确定的路径：
\begin{equation*}
 \Psi(\chi)=\log\sum_vt(v)e^{\chi u_v},\qquad
 q^{[\chi]}(v)=t(v)e^{\chi u_v-\Psi(\chi)}.
\end{equation*}
这里 $\chi$ 是路径坐标，实际第 $k$ 轮取 $\chi=\chi_k$；它不代表新的策略变量。归一化给出 $\Psi(0)=\Psi(1)=0$，以及 $q^{[0]}=t$、$q^{[1]}=q_0$。

\begin{theorem}[KL 路径的精确积分]
\label{dyn:thm:hs-kl-path}
配分函数满足
\begin{equation*}
 \Psi'(\chi)=\E_{q^{[\chi]}}u,
 \qquad\Psi''(\chi)=\Var_{q^{[\chi]}}(u).
\end{equation*}
沿路径的 $\KL(q^{[\chi]}\Vert t)$ 与 $\KL(t\Vert q^{[\chi]})$ 分别为
\begin{align}
 \KL(q^{[\chi]}\Vert t)
 &=\chi\Psi'(\chi)-\Psi(\chi)
 =\int_0^\chi z\Psi''(z)\,\mathrm dz,
 \label{dyn:eq:hs-forward-kl-path}\\
 \KL(t\Vert q^{[\chi]})
 &=\Psi(\chi)-\chi\Psi'(0)
 =\int_0^\chi(\chi-z)\Psi''(z)\,\mathrm dz.
 \label{dyn:eq:hs-reverse-kl-path}
\end{align}
任意路径两点之间还有
\begin{equation*}
 \KL(q^{[\chi]}\Vert q^{[\widetilde\chi]})
 =(\chi-\widetilde\chi)\Psi'(\chi)-\Psi(\chi)+\Psi(\widetilde\chi).
\end{equation*}
\end{theorem}

\begin{proof}
有限词表允许逐项微分。令 $C(\chi)=\sum_vt(v)e^{\chi u_v}$，则
\begin{equation*}
 \Psi'(\chi)=\frac{C'(\chi)}{C(\chi)}=\sum_vq^{[\chi]}(v)u_v,
 \qquad
 \Psi''(\chi)=\frac{C''(\chi)}{C(\chi)}-
          \left(\frac{C'(\chi)}{C(\chi)}\right)^2.
\end{equation*}
右侧正好是 $u$ 的方差。相对 teacher 的 log-ratio 是 $\chi u_v-\Psi(\chi)$，在 $q^{[\chi]}$ 下取期望得到第一式。其关于 $\chi$ 的导数为 $\chi\Psi''(\chi)$，在零点的值为零，积分得到式~\eqref{dyn:eq:hs-forward-kl-path}。相对路径分布的 teacher log-ratio 是 $\Psi(\chi)-\chi u_v$，在 teacher 下取期望得到第二式；再对 $\Psi$ 在零点使用带积分余项的二阶展开，得到式~\eqref{dyn:eq:hs-reverse-kl-path}。最后两条路径分布的 log-ratio 为 $(\chi-\widetilde\chi)u_v-\Psi(\chi)+\Psi(\widetilde\chi)$，取期望即得。
\end{proof}

由于 $\Psi''(\chi)\geq0$，$\KL(q^{[\chi]}\Vert t)$ 在 $\chi>0$ 上随 $\chi$ 增大而增大，在 $\chi<0$ 上随 $\chi$ 增大而减小。$\KL(t\Vert q^{[\chi]})$ 的导数 $\Psi'(\chi)-\Psi'(0)$ 也具有相同符号。因此 $0<\eta\leq1$ 下，两个方向的 KL 都逐轮下降。$1<\eta<2$ 时，连续两轮的 $\chi$ 符号不同，跨度收缩本身不能比较它们的 KL 大小。

\begin{corollary}[局部 KL 速率与全局绝对界]
非 teacher 初值、固定 $0<\eta<2$ 且 $\eta\neq1$ 时，
\begin{equation*}
 \lim_{k\to\infty}
 \frac{\KL(q_{k+1}\Vert t)}{\KL(q_k\Vert t)}
 =(1-\eta)^2.
\end{equation*}
对所有实数 $\chi_k$，还有
\begin{equation*}
 \KL(q_k\Vert t)\leq\frac{\chi_k^2b_0^2}{8}.
\end{equation*}
\end{corollary}

\begin{proof}
在零点展开得到
\begin{equation*}
 \KL(q^{[\chi]}\Vert t)
 =\frac{\chi^2}{2}\Var_{t}(u)+O(|\chi|^3).
\end{equation*}
非 teacher 初值使该方差严格正；将 $\chi=\chi_k$ 与 $\chi_{k+1}=(1-\eta)\chi_k$ 代入并相除，得到极限。任意支撑在长度为 $b_0$ 的区间内的随机变量，方差不超过 $b_0^2/4$。把这个界代入式~\eqref{dyn:eq:hs-forward-kl-path}，得到 $\chi^2b_0^2/8$；当 $\chi<0$ 时反向积分同样给出正的 $\chi^2/2$ 因子。
\end{proof}

这个全局绝对界以初始跨度衡量误差；它没有把当前 KL 与初始 KL 作比例比较。局部极限 $(1-\eta)^2$ 也不能提升为全局单步收缩因子。

\subsection{精确三点恒等式与单步下降}

\begin{theorem}[Hard sync 的三点 KL 恒等式]
\label{dyn:thm:hs-three-point}
记 $q^+=T_\eta(q)$。对任意 $\eta>0$，
\begin{equation}\label{dyn:eq:hs-three-point}
 \eta\KL(q\Vert t)
 =\eta\KL(q^+\Vert t)
 +(1-\eta)\KL(q^+\Vert q)
 +\KL(q\Vert q^+).
\end{equation}
\end{theorem}

\begin{proof}
从更新定义分别减去 anchor 和 teacher 的 log-density：
\begin{align*}
 \log\frac{q^+}{q}
 &=-\eta\log\frac{q}{t}-\log Z_\eta(q),\\
 \log\frac{q^+}{t}
 &=(1-\eta)\log\frac{q}{t}-\log Z_\eta(q).
\end{align*}
第一式在 $q$ 下取负期望，得到
\begin{equation*}
 \KL(q\Vert q^+)
 =\eta\KL(q\Vert t)+\log Z_\eta(q).
\end{equation*}
第二式在 $q^+$ 下取期望，并与第一式在 $q^+$ 下的期望组合，得到
\begin{equation*}
 \eta\KL(q^+\Vert t)
 +(1-\eta)\KL(q^+\Vert q)
 =-\log Z_\eta(q).
\end{equation*}
两式相加，配分函数消去，即为所需恒等式。
\end{proof}

当 $0<\eta\leq1$ 时，式~\eqref{dyn:eq:hs-three-point} 的最后两个项非负。常步长的下降量和累计更新量因此满足
\begin{align*}
 \KL(q_k\Vert t)-\KL(q_{k+1}\Vert t)
 &=\frac{1-\eta}{\eta}\KL(q_{k+1}\Vert q_k)
 +\frac1\eta\KL(q_k\Vert q_{k+1}),\\
 \sum_{k=0}^{\infty}\KL(q_k\Vert q_{k+1})
 &\leq\eta\KL(q_0\Vert t).
\end{align*}
后一式由前式对有限轮数求和、去掉非负项，再令轮数趋于无穷得到。它控制整个 hard-sync 轨迹的累计 KL 移动量。

Gibbs target 对任意正步长仍是一个严格凸目标的唯一最优点。具体地，定义
\begin{equation*}
 F_\eta(p)=\KL(p\Vert t)
 +(\eta^{-1}-1)\KL(p\Vert q).
\end{equation*}
直接代入 $q^+$ 的 log-density，得到
\begin{equation*}
 F_\eta(p)=\eta^{-1}\KL(p\Vert q^+)
 -\eta^{-1}\log Z_\eta(q).
\end{equation*}
所以 $\eta>1$ 时虽然第二个 KL 的系数为负，整个目标的熵系数仍为 $1/\eta>0$，最优点仍唯一。最优点存在、hard-sync 稳定和 teacher KL 单步下降是三个不同结论。

\subsection{KL 收缩的两个数值反例}

取二元词表，$q_0=(10^{-12},1-10^{-12})$、$t=(0.1,0.9)$、$\eta=1.1$。直接代入更新得到
\begin{equation*}
 q_1=(0.5856848855168,\;0.4143151144832),
\end{equation*}
\begin{equation*}
 \KL(q_0\Vert t)=0.1053605156314,
 \qquad\KL(q_1\Vert t)=0.7138510884672.
\end{equation*}
跨度严格缩小到初值的 $0.1$，$\KL(q^{[\chi]}\Vert t)$ 却增加到约 $6.7753$ 倍。这说明 $1<\eta<2$ 的稳定性不能表述为每一步 teacher KL 下降。

取 $q_0=(1-10^{-12},10^{-12})$、$t=(0.01,0.99)$、$\eta=0.5$，得到
\begin{equation*}
 q_1=(0.9999900502246,\;0.0000099497753721),
\end{equation*}
\begin{equation*}
 \KL(q_0\Vert t)=4.6051701859549,
 \qquad\KL(q_1\Vert t)=4.6049999147314.
\end{equation*}
此时 KL 的比值约为 $0.9999630$，超过 $1-\eta=0.5$ 与 $(1-\eta)^2=0.25$。这一次 KL 确实下降，但远处的下降速度不能套用 teacher 附近的二阶渐近速率。

\subsection{从跨度到分布误差的最优 TV 界}

\begin{theorem}[Hilbert 距离控制 TV 的 sharp 界]
\label{dyn:thm:hs-sharp-tv}
任意严格正概率分布 $p,q$ 满足
\begin{equation}\label{dyn:eq:hs-sharp-tv}
 \TV(p,q)\leq\tanh\!\left(\frac{d_H(p,q)}4\right).
\end{equation}
对每个给定正跨度，都存在二元分布达到等号。
\end{theorem}

\begin{proof}
令 $L_v=p(v)/q(v)$、$m=\min L_v$、$M=\max L_v$、$\delta=d_H(p,q)$。归一化给出 $\E_qL=1$，所以 $m\leq1\leq M$，且 $M/m=e^\delta$。$\delta=0$ 时两个分布相等，结论成立。以下设 $\delta>0$。

凸函数 $|L-1|$ 在 $[m,M]$ 上不超过端点连接的弦线，因此
\begin{equation*}
 |L-1|\leq\frac{M-L}{M-m}(1-m)
 +\frac{L-m}{M-m}(M-1).
\end{equation*}
在 $q$ 下取期望，再除以二，得到
\begin{equation*}
 \TV(p,q)=\tfrac12\E_q|L-1|
 \leq\frac{(M-1)(1-m)}{M-m}.
\end{equation*}
记 $h=e^\delta$，代入 $M=hm$ 将右侧化为
\begin{equation*}
 \frac{h+1-hm-m^{-1}}{h-1}.
\end{equation*}
由 $hm+m^{-1}\geq2\sqrt h$，继续得到
\begin{equation*}
 \TV(p,q)\leq\frac{\sqrt h-1}{\sqrt h+1}
 =\tanh(\delta/4).
\end{equation*}
为了达到等号，取两个 token 的概率比分别为 $M=e^{\delta/2}$ 和 $m=e^{-\delta/2}$，在 $q$ 下赋予质量
\begin{equation*}
 q(L=M)=\frac1{1+e^{\delta/2}},\qquad
 q(L=m)=\frac{e^{\delta/2}}{1+e^{\delta/2}}.
\end{equation*}
这些质量满足 $\E_qL=1$。随机变量只取端点值，使弦线界取等号；$hm=m^{-1}$ 又使算术--几何均值界取等号。故整个界达到等号。
\end{proof}

结合 hard-sync 闭式立即得到
\begin{equation*}
 \TV(q_k,t)\leq\tanh(|\chi_k|b_0/4).
\end{equation*}
对于常步长 $0<\eta<2$，若要求 TV 不超过 $\delta\in(0,1)$，充分条件是
\begin{equation*}
 |1-\eta|^k b_0\leq4\operatorname{artanh}(\delta).
\end{equation*}
当 $b_0$ 已不超过右侧时，初始策略就满足要求；否则取 $\eta\neq1$，可解出所需 outer 轮数
\begin{equation*}
 k\geq
 \frac{\log\!\left(b_0/[4\operatorname{artanh}(\delta)]\right)}
 {-\log|1-\eta|}.
\end{equation*}
这个保证以逐 prefix 的精确 Gibbs 更新为前提。它没有把有限 inner、参数共享或 rollout 测度变化隐藏到精确公式中。
